\title{Controlling Text-to-Image Diffusion by Orthogonal Finetuning}

\begin{abstract}
Large-scale text-to-image diffusion models demonstrate remarkable proficiency in producing photorealistic imagery based on textual descriptions. A significant unresolved question concerns how to efficiently steer or adapt these robust models to accommodate a range of downstream applications. In response to this, we propose Orthogonal Finetuning~(OFT), a systematic finetuning strategy for adjusting text-to-image diffusion models for specific downstream requirements. Distinct from prior approaches, OFT is theoretically guaranteed to maintain hyperspherical energy, an attribute that reflects the relationships between neurons on the unit hypersphere. Our analysis indicates that conserving this feature is vital for retaining the original semantic generation capabilities of text-to-image diffusion models. To further promote stable finetuning, we introduce Constrained Orthogonal Finetuning (COFT), which enforces an added constraint on the hyperspherical radius. In our study, we focus on two prominent finetuning applications: subject-driven generation—where, given several images of a particular subject and a text prompt, the objective is to synthesize novel images of that subject in new settings—and controllable generation, where the aim is to augment the model with additional control signals. Experimental evidence demonstrates that our OFT methodology yields superior generation quality and accelerates convergence relative to current techniques.
\end{abstract}

\section{Introduction}

Recent advances in text-to-image diffusion models~\cite{saharia2022photorealistic,ramesh2022hierarchical,rombach2022high} have set a new benchmark for high-quality, text-conditioned image synthesis. However, the degree of control achievable through text input alone often falls short, leaving fine-grained and precise manipulation of generated outputs as an open issue. To overcome these limitations, we address two classes of text-to-image synthesis problems in this work:

\begin{itemize}[leftmargin=*,nosep]

    \item \textbf{Subject-driven generation}~\cite{ruiz2023dreambooth}: The task involves, given a small set of example images of a particular subject, creating new images depicting the same subject in diverse scenarios as specified by a text prompt.
    \item \textbf{Controllable generation}~\cite{zhang2023adding,mou2023t2i}: Here, the objective is to enable image synthesis based on supplemental control inputs (such as canny edge maps or segmentation masks) in conjunction with a textual description.
\end{itemize}

At their core, both tasks center on devising methods to adapt text-to-image diffusion models to new objectives without compromising the generative quality established during pretraining. The primary requirements for an optimal finetuning approach can be articulated as follows: (1) \emph{training efficiency}, reflected by minimizing both the number of learnable parameters and training epochs, and (2) \emph{generalizability preservation}, which entails maintaining the diversity and fidelity of generated outputs. Common finetuning methodologies involve either updating neuron weights incrementally with a low learning rate (\emph{e.g.}, \cite{ruiz2023dreambooth}), or integrating a compact module parameterized with its own set of weights (\emph{e.g.}, \cite{hulora2022,zhang2023adding}). However, these conventional strategies lack guarantees that the original generative performance will remain intact. Moreover, there is an absence of systematic frameworks for crafting effective finetuning procedures or selecting hyperparameters such as the number of epochs. One significant challenge stems from the unavailability of a robust metric to assess the retention of pretrained generative capabilities. Existing approaches tend to operate under the assumption that a minimal Euclidean distance between the adjusted and original model weights corresponds to better preservation of the pretrained attributes. Consequently, finetuning is generally conducted with extremely small learning rates. Nonetheless, this underlying premise is not universally valid, as the Euclidean distance alone fails to adequately reflect the degree of semantic preservation. Thus, we propose that the adoption of a more structurally informed metric to evaluate the deviation between the finetuned and pretrained models could enhance both the stability of the finetuning process and the fidelity of generative performance inherited from pretraining.

Drawing upon empirical findings that hyperspherical similarity effectively represents semantic relationships~\cite{liu2017deep,liu2018decoupled,chen2020angular}, we utilize the concept of hyperspherical energy~\cite{liu2018learning} to describe the pairwise relational structure among neurons. Hyperspherical energy, quantified as the aggregated hyperspherical similarity (\emph{e.g.}, cosine similarity) between all neuron pairs in a given layer, reflects the degree of uniformity in neuron distribution over the unit hypersphere~\cite{liu2021learning}. Our working hypothesis is that an optimal finetuned model should display minimal deviation in hyperspherical energy relative to its pretrained counterpart. A straightforward approach would be to introduce a regularization term that constrains hyperspherical energy during finetuning; however, this method does not necessarily minimize the energy difference effectively. Instead, we leverage a key invariance property: applying a uniform orthogonal transformation to all neurons preserves their pairwise hyperspherical similarities. Based on this invariance, we put forward Orthogonal Finetuning~(OFT), a method for adapting large text-to-image diffusion models to downstream tasks while maintaining invariant hyperspherical energy. The principal mechanism is to learn an orthogonal transformation, shared across layers, that preserves inter-neuron angular relationships. Alternatively, OFT can be interpreted as modifying the underlying coordinate framework for neurons within the same layer. To further enhance pretraining performance retention, we propose Constrained Orthogonal Finetuning~(COFT), a variant that confines the finetuned model to reside on a hypersphere of fixed radius centered at the pretrained neuron locations, guided by the observation that minimizing the Euclidean distance between pretrained and finetuned models better conserves pretraining characteristics.

The rationale behind employing orthogonal transformation for finetuning neurons is partly inspired by the 2D Fourier transform, which allows an image to be expressed through its magnitude and phase spectra. Notably, the phase spectrum—representing angular relationships between inputs and basis elements—captures the essential semantic content. Empirically, retaining the phase spectrum and replacing the magnitude with random values often suffices to reconstruct the original image’s semantic features. This observation underlines that adjusting the orientation of neurons primarily influences the semantic aspects of the generated output, aligning with the objectives of both subject-driven and controllable generation. However, unconstrained manipulation of neuron directions can compromise the generative abilities established during pretraining. To prevent this, we advocate for maintaining the pairwise angles between neurons, premised on the notion that these angular relationships are fundamental to encoding network knowledge. Accordingly, we propose a layer-wise, shared orthogonal transformation for neurons, designed so that the hyperspherical energy remains constant.

Additionally, our approach draws on orthogonal over-parameterized training~\cite{liu2021orthogonal}, where neural networks for classification are initialized randomly and then trained by applying orthogonal transformations. In this method, random initializations yield a low expected hyperspherical energy, and the strategy is to maintain this low energy throughout training—a property linked to improved generalization in classification tasks~\cite{liu2018learning,lin2020regularizing}. The work in~\cite{liu2021orthogonal} demonstrates the sufficiency of orthogonal transformations for constructing neural networks that generalize well in classification settings. In contrast, our focus lies in adapting text-to-image diffusion models for enhanced controllability and robust generative capabilities. The distinctions between OFT and~\cite{liu2021orthogonal} are twofold. Firstly, whereas~\cite{liu2021orthogonal} seeks to minimize hyperspherical energy, OFT’s objective is to preserve the hyperspherical energy of the original pretrained model, ensuring that the structural properties developed during pretraining are not lost during finetuning—a crucial concern for diffusion models where reducing hyperspherical energy may disrupt intrinsic semantic representations. Secondly, OFT emphasizes minimizing deviations from the pretrained network, motivating a constrained approach, in contrast to the unconstrained optimization adopted in~\cite{liu2021orthogonal}. We posit that effective finetuning requires a balance between adaptability and preservation of pretrained structure, a balance our OFT methodology is designed to provide. Below, we outline our specific contributions:

\begin{itemize}[leftmargin=*,nosep]

    \item We introduce a new finetuning approach, Orthogonal Finetuning (OFT), designed to enhance the controllability of text-to-image diffusion models. To promote additional training stability, we also devise a constrained variant that restricts the angle of deviation from the original, pretrained model parameters.
    \item Unlike current finetuning techniques, OFT enables parameter adaptation while formally maintaining hyperspherical energy—an attribute we empirically recognize as crucial for upholding the semantic generative fidelity of the pretrained model.
    \item We employ OFT in two scenarios: subject-driven image generation and controlled generation tasks. Through extensive experiments, we observe that OFT consistently outperforms existing methods regarding image generation quality, training convergence rate, and the robustness of finetuning. Additionally, OFT demonstrates higher sample efficiency, achieving effective convergence using significantly fewer training images and epochs.
    \item For evaluating controllability, we present a novel metric termed control consistency. This method assesses the match between control signals inferred from generated images and their original control inputs. Results using this metric further substantiate OFT’s strong capabilities for controlled generation.
\end{itemize}

\section{Related Work}

\textbf{Text-to-image diffusion models}. The field of text-to-image synthesis has witnessed substantial advancements~\cite{saharia2022photorealistic,ramesh2022hierarchical,rombach2022high,gu2022vector,nichol2022glide}, which are largely attributable to progress in diffusion-based generative frameworks~\cite{ho2020denoising,song2021score,song2021denoising,dhariwal2021diffusion} alongside innovations in vision-language learning~\cite{su2020vlbert,lu2019vilbert,radford2021learning,wang2021simvlm,li2022blip,li2023blip,alayrac2022flamingo,schuhmann2022laion}. In approaches such as GLIDE~\cite{nichol2022glide} and Imagen~\cite{saharia2022photorealistic}, diffusion models operate within the pixel domain. Specifically, GLIDE couples the training of a text encoder with the diffusion model using text-image pairs, whereas Imagen utilizes a fixed, pretrained text encoder throughout. Conversely, techniques like Stable Diffusion~\cite{rombach2022high} and DALL-E2~\cite{ramesh2022hierarchical} perform diffusion in latent spaces. Stable Diffusion relies on a visual codebook learned via VQ-GAN~\cite{esser2021taming} to construct its latent representations, while DALL-E2 makes use of CLIP~\cite{radford2021learning} to jointly embed images and text into a shared latent space. Beyond diffusion models, other architectures such as generative adversarial networks~\cite{reed2016generative,zhang2017stackgan,xu2018attngan,li2019controllable} and autoregressive models~\cite{ramesh2021zero,ding2021cogview,wu2022nuwa,yuscaling} have also been applied to the text-to-image task. OFT functions independently of specific model architecture and thus can be incorporated with any text-to-image diffusion system.

\textbf{Subject-driven generation}. To reduce unwanted changes to the subject, prior approaches~\cite{avrahami2022blended,nichol2022glide} utilize user-provided masks as additional input conditions. Inversion strategies~\cite{choi2021ilvr,dhariwal2021diffusion,ramesh2022hierarchical,gal2022image} aim to modify contextual details while keeping the subject intact. The method of~\cite{hertz2022prompt} offers both local and global editing capabilities without dependence on input masks. However, such approaches tend to fall short in retaining fine, identity-specific features. In contrast, Pivotal Tuning~\cite{roich2022pivotal} involves locally updating a generator around an inverted latent code while introducing regularization to maintain identity preservation. Likewise,~\cite{nitzan2022mystyle} builds a personalized face generator from a set of individual face images. Although~\cite{casanova2021instance} is capable of producing diverse samples from a single instance, it may neglect distinctive instance details. DreamBooth~\cite{ruiz2023dreambooth} enables personalizing a model using a new token and a few reference subject images, employing a combination of reconstruction and class-prior preservation losses. Distinguishing itself from DreamBooth's standard finetuning via low learning rates, OFT leverages orthogonal transformations as the mechanism for parameter updates

\textbf{Controllable generation}. Image-to-image translation can be framed as a problem of controllable generation. Earlier work in this area predominantly employs conditional generative adversarial networks~\cite{isola2017image,zhu2017toward,wang2018high,park2019semantic,choi2018stargan}. In addition, diffusion models have gained traction as effective tools for image-to-image translation~\cite{saharia2022palette,voynov2022sketch,wang2022pretraining}. The introduction of ControlNet~\cite{zhang2023adding} marked a significant advancement, allowing fine-grained manipulation of a pretrained diffusion model via finetuning with auxiliary control signals, resulting in strong controllability. Similarly, T2I-Adapter~\cite{mou2023t2i} independently proposes to finetune diffusion models to enhance control over image outputs. Building on these developments~\cite{zhang2023adding,mou2023t2i}, our approach leverages OFT to finetune pretrained diffusion models, achieving superior controllability with less training data and fewer finetuned parameters. Notably, OFT brings no additional computational cost during inference.

\textbf{Model finetuning}. The practice of adapting large pretrained models to downstream tasks through finetuning has become increasingly prevalent~\cite{devlin2018bert,brown2020language,he2022masked}. Within this context, adaptation techniques (see, \emph{e.g.}, \cite{hulora2022,houlsby2019parameter,pfeiffer2020adapterfusion}) are widely investigated in the field of natural language processing. Of these, LoRA~\cite{hulora2022} is most closely related to OFT; it posits that updates during finetuning can be well-approximated by a low-rank additive structure. In contrast, OFT modifies neuron weights via a layer-wise shared orthogonal transformation in a multiplicative fashion, and mathematically guarantees preservation of the pairwise angles between neurons within a layer. This leads to greater training stability.

\section{Orthogonal Finetuning}

\subsection{Why Does Orthogonal Transformation Make Sense?}

To motivate the use of orthogonal transformations in the finetuning of text-to-image diffusion models, we address two interrelated sub-questions: (1) the rationale for adjusting the angle (i.e., direction) of neuron weights, and (2) the justification for employing orthogonal transformations for this purpose.

To address the first question, we build upon empirical findings from prior work \cite{liu2018decoupled,chen2020angular}, which suggest that differences in the angular features of neurons are indicative of the semantic gap. SphereNet~\cite{liu2017deep} demonstrates that constraining all neurons of a neural network to lie on a unit hypersphere not only preserves model capacity but also enhances generalization, indicating that neuronal orientation suffices to encapsulate essential data characteristics. To illustrate the relevance of neuron angles, we design a simple experiment, depicted in Figure~\ref{fig:toy_ae}, by training a standard convolutional autoencoder with flower images. During training, feature maps are computed via the conventional inner product ($z$ is the output of the convolutional filter $\bm{w}$ applied to the input $\bm{x}$ within a sliding window). At test time, we evaluate three methods for generating feature maps: (a) the training inner product, (b) only the norm information, and (c) only the angular information. The outcomes, shown in Figure~\ref{fig:toy_ae}, reveal that the angular features alone almost completely reconstruct the input images, whereas magnitude information is insufficient. Importantly, cosine-based (angular) activation (c) is not introduced during training, which proceeds solely with the inner product. These findings suggest that neuron orientation is chiefly responsible for encoding the semantic content of input images, and adjusting neuron directions should be an effective approach to modify such semantics.

Turning to the second question, a straightforward method to adjust neuron orientation is by collectively rotating or reflecting all neurons within the same layer, a process naturally formalized as an orthogonal transformation. While alternative angular transformations allowing different rotational adjustments per neuron could offer greater flexibility, our experiments indicate that orthogonal transformations provide an optimal compromise between adaptability and structure. Furthermore, results from \cite{liu2021orthogonal} establish that orthogonal transformations possess sufficient expressive power for training neural networks. To empirically validate our approach, we conduct a controllable generation experiment under the ControlNet~\cite{zhang2023adding} protocol, as shown in Figure~\ref{fig:ortho}. Our conventional OFT, under the orthogonality constraint, consistently exhibits stable performance and achieves precise control after twenty epochs of training. By contrast, OFT implemented without the orthogonality condition fails to yield realistic images or control over outputs. This evidence underscores the necessity of enforcing the orthogonality constraint in OFT.

\subsection{General Framework}

Finetuning conventionally relies on gradient descent with a reduced learning rate to make incremental adjustments to a model or selected layers thereof. By employing a low learning rate, this approach inherently discourages significant departures from the pretrained parameters, thus standard finetuning operates by implicitly reducing the difference $\thickmuskip=2mu \medmuskip=2mu \| \bm{M} - \bm{M}^0\|$, with $\bm{M}$ representing the finetuned model parameters and $\bm{M}^0$ denoting the original pretrained weights. While this implicit proximity constraint is intuitively reasonable, it may allow excessive flexibility, particularly in the context of large-scale models. To mitigate this, LoRA augments the update process with an explicit low-rank restriction on weight changes, i.e., $\thickmuskip=2mu \medmuskip=2mu \text{rank}(\bm{M} - \bm{M}^0)=r'$, where $r'$ is deliberately chosen to be small. OFT, contrastingly, imposes a different regularization by requiring the pairwise similarity among neurons to be preserved: $\thickmuskip=2mu \medmuskip=2mu \| \text{HE}(\bm{M})-\text{HE}(\bm{M}^0) \|=0$, where $\text{HE}(\cdot)$ is the hyperspherical energy associated with a given weight matrix.

To provide concrete intuition, consider a dense (fully connected) layer expressed as $\thickmuskip=2mu \medmuskip=2mu \bm{W}=\{\bm{w}_1,\cdots,\bm{w}_n\}\in\mathbb{R}^{d\times n}$ with each $\bm{w}_i\in\mathbb{R}^{d}$ corresponding to the $i$-th neuron (and $\bm{W}^0$ as its pretrained counterpart). For an input vector $\thickmuskip=2mu \medmuskip=2mu \bm{x}\in\mathbb{R}^d$, the output is calculated as $\thickmuskip=2mu \medmuskip=2mu \bm{z}=\bm{W}^\top\bm{x}$ with $\thickmuskip=2mu \medmuskip=2mu \bm{z}\in\mathbb{R}^n$. The OFT framework can thus be interpreted as striving to align the hyperspherical energy of the finetuned weights with that of the pretrained ones:

\begin{equation}\label{eq:oft_principle}
\footnotesize
\min_{\bm{W}} \left\| \text{HE}(\bm{W})-\text{HE}(\bm{W}^0)\right\|~~\Leftrightarrow~~\min_{\bm{W}} \bigg{\|} \sum_{i\neq j} \|\hat{\bm{w}}_i-\hat{\bm{w}}_j\|^{-1}-\sum_{i\neq j} \|\hat{\bm{w}}^0_i-\hat{\bm{w}}^0_j\|^{-1}\bigg{\|}
\end{equation}

Here, each $\thickmuskip=2mu \medmuskip=2mu \hat{\bm{w}}_i:=\bm{w}_i/\|\bm{w}_i\|$ is a unit vector representing the $i$-th neuron, and the hyperspherical energy for a layer $\bm{W}$ is given by $\thickmuskip=2mu \medmuskip=2mu \text{HE}(\bm{W}):= \sum_{i\neq j} \|\hat{\bm{w}}_i-\hat{\bm{w}}_j\|^{-1}$. Notably, the lower bound of the objective in Eq.~\eqref{eq:oft_principle} is zero, attainable whenever $\bm{W}$ and $\bm{W}^0$ are connected by an orthogonal transformation, specifically when $\thickmuskip=2mu \medmuskip=2mu \bm{W}=\bm{R}\bm{W}^0$ with $\thickmuskip=2mu \medmuskip=2mu \bm{R}\in\mathbb{R}^{d\times d}$ orthogonal (determinant $1$ for rotation, $-1$ for reflection). The fundamental insight behind OFT is that it restricts fin

\begin{equation}\label{eq:oft_general}
    \footnotesize
    \bm{z}=\bm{W}^\top\bm{x}=(\bm{R}\cdot\bm{W}^0)^\top\bm{x},~~~\text{s.t.}~\bm{R}^\top\bm{R}=\bm{R}\bm{R}^\top=\bm{I}
\end{equation}

Here, $\bm{W}$ refers to the weight matrix obtained through OFT finetuning, while $\bm{I}$ stands for the identity matrix. The structure of OFT is presented in Figure~\ref{fig:block}. Analogous to LoRA's approach of zero initialization, it is crucial that OFT allows finetuning to begin precisely from the pretrained weights $\bm{W}^0$. To realize this, we initialize the orthogonal matrix $\bm{R}$ as the identity matrix, ensuring the starting point of finetuning is aligned with the original pretrained parameters. The orthogonality constraint on $\bm{R}$ can be enforced using differentiable orthogonalization methods as described in \cite{liu2021orthogonal,lezcano2019cheap}. Later, we discuss efficient means to maintain this orthogonality constraint.

\subsection{Efficient Orthogonal Parameterization}\label{sect:eff_ortho}

Although classical techniques like the Gram-Schmidt process yield differentiable orthogonalization, their computational demands make them impractical for large-scale applications~\cite{liu2021orthogonal}. For enhanced computational efficiency, we leverage the Cayley parameterization to generate orthogonal matrices. More precisely, we obtain the orthogonal matrix by the formula $\thickmuskip=2mu \medmuskip=2mu \bm{R}=(\bm{I}+\bm{Q})(I-\bm{Q})^{-1}$, where $\bm{Q}$ denotes a skew-symmetric matrix, i.e., $\thickmuskip=2mu \medmuskip=2mu \bm{Q}=-\bm{Q}^\top$. The main trade-off of using the Cayley parameterization is that it can only construct orthogonal matrices with determinant $1$, restricting $\bm{R}$ to the special orthogonal group. Nevertheless, this limitation has not been observed to negatively impact practical performance. Still, Cayley parameterized $\bm{R}$ remains potentially parameter-inefficient for large $d$. To address this, we design $\bm{R}$ as a block-diagonal matrix composed of $r$ blocks, giving:

\begin{equation}
    \footnotesize
    \bm{R}=\text{diag}(\bm{R}_1,\bm{R}_2,\cdots,\bm{R}_r)=\begin{bmatrix}
    \bm{R}_1\in \text{O}(\frac{d}{r}) &  &\\
    &\ddots & \\
    & & \bm{R}_r\in \text{O}(\frac{d}{r})
    \end{bmatrix}\in \text{O}(d)
\end{equation}

Here, $\text{O}(d)$ designates the orthogonal group in $d$-dimensional space, where $\thickmuskip=2mu \medmuskip=2mu \bm{R}\in\mathbb{R}^{d\times d}$ and each $\thickmuskip=2mu \medmuskip=2mu \bm{R}_i\in\mathbb{R}^{d/r\times d/r}$ for all $i$ are orthogonal matrices. In the special case of $\thickmuskip=2mu \medmuskip=2mu r=1$, the block-diagonal orthogonal matrix reduces to a conventional, unrestricted orthogonal matrix. The parameter count for a $d\times d$ orthogonal matrix is given by $\thickmuskip=2mu \medmuskip=2mu d(d-1)/2$, which leads to a computational complexity of $\mathcal{O}(d^2)$. Conversely, for an $r$-block diagonal orthogonal matrix, the parameter count becomes $\thickmuskip=2mu \medmuskip=2mu d(d/r-1)/2$, corresponding to a complexity of $\thickmuskip=2mu \medmuskip=2mu \mathcal{O}(d^2/r)$. If further parameter reduction is desired, one may impose sharing among the block matrices—specifically, setting $\thickmuskip=2mu \medmuskip=2mu\bm{R}_i=\bm{R}_j$ for all $i\neq j$—which decreases the parameter complexity to $\mathcal{O}(d^2/r^2)$. Importantly, these approaches do not compromise the orthogonality of $\bm{R}$, ensuring that hyperspherical energy is preserved.

We further analyze OFT's parameter efficiency relative to LoRA. For LoRA, using a low-rank decomposition with rank $r'$, the number of trainable parameters is $\thickmuskip=2mu \medmuskip=2mu r'(d+n)$. If both $r$ (in OFT) and $r'$ (in LoRA) are chosen as fractions of the neuron dimension $d$—that is, $\thickmuskip=2mu \medmuskip=2mu r=r'=\alpha d$ for some $\thickmuskip=2mu \medmuskip=2mu 0<\alpha\leq 1$—then LoRA's parameter complexity scales as $\mathcal{O}(d^2+dn)$, whereas for OFT the scaling is $\mathcal{O}(d)$. For instance, consider a weight matrix of dimensions $\thickmuskip=2mu \medmuskip=2mu 128\times 128$: LoRA with $\thickmuskip=2mu \medmuskip=2mu r'=8$ involves $2,048$ learnable parameters, whereas OFT with $\thickmuskip=2mu \medmuskip=2mu r=8$ (without block sharing) only requires $960$.

\subsection{Constrained Orthogonal Finetuning}

The flexibility inherent to standard OFT may be further curtailed by restricting the finetuned parameters to remain within a proximity of the initial, pretrained model. In practice, Constrained Orthogonal Finetuning (COFT) executes the following transformation in the forward computation:

\begin{equation}\label{eq:coft_general}
    \footnotesize
    \bm{z}=\bm{W}^\top\bm{x}=(\bm{R}\cdot\bm{W}^0)^\top\bm{x},~~~\text{s.t.}~\bm{R}^\top\bm{R}=\bm{R}\bm{R}^\top=\bm{I},~\norm{\bm{R}-\bm{I}}\leq \epsilon
\end{equation}

Here, an orthogonality requirement on $\bm{R}$ is paired with an $\epsilon$-proximity constraint to the identity matrix. While orthogonality may be imposed by employing the Cayley parameterization (see Section~\ref{sect:eff_ortho}), enforcing the $\epsilon$-deviation constraint within the Cayley framework poses additional challenges. To better understand the effect of the Cayley transform, we utilize the Neumann series to approximate $\thickmuskip=2mu \medmuskip=2mu \bm{R}=(\bm{I}+\bm{Q})(I-\bm{Q})

Since the series converges in the operator norm, the constraint $\thickmuskip=2mu \medmuskip=2mu\|\bm{R}-\bm{I}\|\leq\epsilon$ can be transferred inside the Cayley transform, resulting in an equivalent condition $\thickmuskip=2mu \medmuskip=2mu \|\bm{Q}-\bm{0}\|\leq\epsilon'$. Here, $\bm{0}$ stands for the zero matrix, and $\epsilon'$ is a different, independent hyperparameter that controls the approximation error. This updated restriction on $\bm{Q}$ can be readily implemented through projected gradient descent. To initialize the orthogonal matrix $\bm{R}$ as the identity, we simply start with $\bm{Q}$ set to zero. The COFT approach integrates two concrete constraints: limiting hyperspherical energy deviation and explicitly restricting divergence from the pretrained parameters. While most finetuning techniques incorporate the latter implicitly, COFT makes this deviation constraint explicit. We find that, despite OFT’s strong empirical results, COFT enhances finetuning stability even further by explicitly bounding deviation. Figure~\ref{fig:eps_coft} illustrates the effect of varying $\epsilon$ on COFT’s performance: as $\epsilon$ increases, the resulting COFT-finetuned models behave increasingly like those trained with OFT, while lower values of $\epsilon$ make COFT’s behavior closely resemble that of the original pretrained text-to-image diffusion network.

\subsection{Re-scaled Orthogonal Finetuning}

As a straightforward augmentation to OFT, we introduce a learnable scaling factor for each neuron. This extension draws from the observation that scaling the neurons' output has no impact on hyperspherical energy due to normalization. Concretely, the forward computation is formulated as: $\thickmuskip=2mu \medmuskip=2mu\bm{z}=(\bm{R}\bm{W}^0\bm{D})^\top\bm{x}$\footnote{\textbf{Errata}: The NeurIPS camera-ready paper inadvertently writes the forward pass for re-scaled OFT as $\thickmuskip=2mu \medmuskip=2mu\bm{z}=(\bm{D}\bm{R}\bm{W}^0)^\top\bm{x}$; however, the actual code implementation matches what is shown here and thus results in Appendix~\ref{app:rescaled} remain accurate.}, where $\thickmuskip=2mu \medmuskip=2mu \bm{D}=\text{diag}(s_1,\cdots,s_n)\in\mathbb{R}^{n\times n}$ is a learnable diagonal matrix with strictly positive entries $s_1,\cdots,s_n$. Unlike the original OFT formulation in Eq.~\eqref{eq:oft_general}, which only updates $\bm{R}$, the re-scaled version jointly optimizes both $\bm{R}$ and the diagonal scaling matrix $\bm{D}$. This modification gives OFT added expressivity with a minimal increase in parameter count. For our main experimental results, we retain the standard OFT to isolate the effect of orthogonal transformation, but our findings indicate that re-scaled OFT generally performs better (see Appendix~\ref{app:rescaled}).

\section{Intriguing Insights and Discussions}

\textbf{OFT is architecture-agnostic}. In theory, OFT can be seamlessly incorporated into any neural network structure. In Transformer models, LoRA is most commonly used for attention weights~\cite{hulora2022}, so for a fair comparison, our experiments also restrict OFT to attention parameters. Beyond dense (fully connected) layers, OFT is a strong candidate for tuning convolutional layers: the block-diagonal form of $\bm{R}$ offers specific interpretations in the context of convolutions, in contrast to LoRA. For instance, if we select the number of blocks to match the number of input channels, each block operates exclusively on a single neuron channel—mirroring the effect of learning depth-wise convolution filters~\cite{chollet2017xception}. Alternatively, if we share the same block in $\bm{R}$ across all channels, each neuron's transformation is governed by a channel-shared orthogonal matrix. We provide initial experiments investigating the use of OFT for convolutional layers in Appendix~\ref{app:convolution}.

\textbf{Connection to LoRA}. LoRA maintains the integrity of the pretrained weight matrix by introducing a low-rank additive component, which restricts major deviations from its original state. In comparison, OFT applies a trainable orthogonal (full-rank) transformation to the pretrained matrix, thereby preserving its preexisting information by constraining the nature of the transformation. The OFT forward computation can be represented as $\thickmuskip=2mu \medmuskip=2mu \bm{z}=(\bm{R}\bm{W}^0)^\top\bm{x}=(\bm{W}^0+(\bm{R}-\bm{I})\bm{W}^0)^\top\bm{x}$, where the term $\thickmuskip=2mu \medmuskip=2mu (\bm{R}-\bm{I})\bm{W}^0$ serves a similar role to LoRA’s low-rank update. Since $\bm{W}^0$ commonly possesses full rank, the modification introduced by OFT corresponds to a low-rank update when $\thickmuskip=2mu \medmuskip=2mu \bm{R}-\bm{I}$ itself is low-rank. Analogous to LoRA’s tunable rank hyperparameter $r'$, OFT employs a block-diagonal parameter $r$ to decrease the trainable parameter count. Fundamentally, OFT and LoRA exemplify two independent strategies for parameter efficiency: LoRA leverages low-rank approximations, whereas OFT reduces parameterization by enforcing block-diagonal sparsity and orthogonality.

\textbf{Why OFT converges faster}. The visual evidence in Figure~\ref{fig:toy_ae} highlights that shifting neuron orientations is the most impactful means of altering model semantics, precisely aligning with OFT’s intended function. Further, OFT may be interpreted as adjusting neurons constrained to a smooth hypersphere, fostering a more favorable optimization surface. Such an approach has been empirically supported in \cite{liu2021orthogonal}.

\textbf{Why not minimize hyperspherical energy}. Unlike the approach in \cite{liu2021orthogonal}, our objective is not to achieve minimum hyperspherical energy. In classification contexts, optimal performance calls for neurons to be non-redundant, but minimizing hyperspherical energy simply results in equidistant neuron spacing on the hypersphere—an outcome that can disrupt pretrained knowledge and is thus unsuitable for finetuning tasks.

\textbf{Balancing flexibility and regularization in finetuning}. Our investigation reveals an inherent tension between achieving flexibility and maintaining regularization. While conventional finetuning methods offer maximum adaptability, they often compromise stability and can easily result in model degeneration. OFT, despite its conceptual simplicity, strikes a favorable equilibrium by ensuring the preservation of pairwise neuron angular relationships. The block-diagonal structure acts as an even stricter form of regularization for the orthogonal transformation.

\textbf{Zero inference overhead}. In contrast to ControlNet, our OFT method does not incur any extra inference costs for the adapted model. During inference, the learned orthogonal transformation $\bm{R}$ can be directly applied to the pretrained weights $\bm{W}^0$, resulting in an updated weight matrix $\thickmuskip=2mu \medmuskip=2mu \bm{W}=\bm{R}\bm{W}^0$ that preserves computational efficiency identical to the original pretrained model.

\section{Experiments and Results}

\textbf{Experimental configurations}. We conduct our experiments using Stable Diffusion v1.5~\cite{rombach2022high} as the foundation for pretrained text-to-image synthesis. To maintain objectivity, images generated by each approach are selected at random. For tasks involving subject-driven image creation, our protocol aligns with DreamBooth~\cite{ruiz2023dreambooth}; for controllable image synthesis, we draw methodology from ControlNet~\cite{zhang2023adding} and T2I-Adapter~\cite{mou2023t2i}. To facilitate equitable comparison with LoRA, we restrict the application of OFT or COFT to the identical model layers where LoRA is employed. Additional implementation details and extensive results appear in Appendix~\ref{app:settings}.

\subsection{Subject-driven Generation}

\textbf{Experimental setup}. We benchmark against DreamBooth~\cite{ruiz2023dreambooth} and LoRA~\cite{hulora2022}. Every approach utilizes the same loss formulation as described in DreamBooth. Both DreamBooth and LoRA are run according to their respective original protocols and are configured with their optimal hyperparameters. Supplementary experimental outcomes can be found in Appendix~\ref{app:settings},\ref{app:coft_oft},\ref{app:more_qualitative},\ref{app:failure}.

\textbf{Finetuning stability and convergence}. We begin by assessing the stability and rate of convergence during finetuning for DreamBooth, LoRA, OFT, and COFT. The corresponding experimental results are displayed in Figure~\ref{fig:teaser} and Figure~\ref{fig:db_convergence}. It is evident that COFT and OFT demonstrate stable finetuning behavior for the diffusion model. After 400 training steps, both DreamBooth and the OFT-based approaches effectively maintain subject control, whereas LoRA struggles to retain subject identity. As training proceeds to 2000 steps, DreamBooth begins producing degenerate outputs, and LoRA fails to depict the yellow shirt, generating yellow fur instead. By contrast, OFT and COFT continue to exhibit robust and reliable control over the generated results. These findings substantiate that our OFT approach achieves rapid and stable convergence in subject-driven generative tasks. Notably, the robustness with respect to the number of finetuning iterations is of practical significance, as it reduces the need for per-subject iteration tuning. For both OFT and COFT, one can conveniently specify a larger iteration count without extensive tuning. Furthermore, when COFT is used with an appropriate $\epsilon$, both the learning rate and the iteration count are largely insensitive, making parameter selection straightforward.

\textbf{Quantitative comparison}. Adopting the evaluation procedure from \cite{ruiz2023dreambooth}, we perform a quantitative analysis of subject fidelity (using DINO~\cite{caron2021emerging} and CLIP-I~\cite{radford2021learning}), fidelity to textual prompts (CLIP-T~\cite{radford2021learning}), as well as diversity among generated samples (LPIPS~\cite{zhang2018unreasonable}). CLIP-I calculates the mean pairwise cosine similarity between CLIP representations for generated versus real images, while DINO operates similarly but employs ViT S/16 DINO features. CLIP-T assesses the mean cosine similarity between textual prompt and the generated images' embeddings. Additionally, sample diversity is gauged by measuring the mean LPIPS cosine similarity among images produced for the same subject and prompt. According to Table~\ref{table:dreambooth_final}, both COFT and OFT surpass DreamBooth and LoRA by a substantial margin in DINO and CLIP-I subject fidelity measures, while also achieving equal or marginally improved results in text prompt alignment and diversity. For reliability, each method is evaluated by finetuning the same pretrained model 30 times with distinct random seeds, thereby reducing variability due to randomness. Collectively, these outcomes confirm that the OFT framework offers superior convergence properties, enhanced stability, and consistently improved performance.

\textbf{Qualitative comparison}. To provide a clearer, more intuitive illustration of the advantages of OFT, we present several randomly selected instances of subject-driven generation in Figure~\ref{fig:dreambooth_final}. For objectivity, all methods generate results using the identical finetuned model, without employing any individually optimized text prompts for the final outputs. Within each method, we choose the version that yields the top-performing validation CLIP scores. As seen in Figure~\ref{fig:dreambooth_final}, both OFT and COFT consistently achieve strong semantic preservation of the subject, while LoRA frequently struggles with maintaining subject identity (\emph{e.g.}, in the bowl example, LoRA completely loses the subject). Additionally, OFT and COFT exhibit superior ability in aligning image generation with the provided text prompt; by contrast, although DreamBooth tends to preserve subject identity, it often fails to follow the prompt's instructions (\emph{e.g.}, the first row of the bowl example). Overall, this comparison highlights that the OFT approach excels in balancing both image controllability and subject fidelity. Furthermore, OFT's effectiveness remains robust across a wide range of iteration counts, unlike DreamBooth and LoRA, whose performance tends to degrade with increased iterations. Additional qualitative samples can be found in Appendix~\ref{app:more_qualitative}. A human study detailed in Appendix~\ref{app:human} offers further empirical support for OFT's superiority.

\subsection{Controllable Generation}

\textbf{Settings}. As comparative baselines, we include ControlNet~\cite{zhang2023adding}, T2I-Adapter~\cite{mou2023t2i}, and LoRA~\cite{hulora2022}. In our main experiments, we tackle three particularly demanding tasks in controllable generation: Canny edge to image~(C2I) using the COCO dataset~\cite{lin2014microsoft}, segmentation map to image~(S2I) with ADE20K~\cite{zhou2017scene}, and landmark to face~(L2F) utilizing the CelebA-HQ dataset~\cite{karras2018progressive,xia2021tedigan}. All methods are finetuned on Stable Diffusion~(SD) v1.5 across each dataset for a total of 20 epochs. Expanded results can be found in Appendix~\ref{app:more_qualitative},\ref{app:more_control_task},\ref{app:failure}.

\textbf{Convergence}. To assess how rapidly ControlNet, T2I-Adapter, LoRA, and COFT converge on the L2F task, we carry out both quantitative and qualitative analyses. As a quantitative metric, we calculate the mean $\ell_2$ distance between the control face landmarks and the landmarks generated by each model. Figure~\ref{fig:face_convergence} displays the progression of face landmark error across different epochs of model fine-tuning. It is evident that COFT and OFT substantially outpace the alternatives in terms of convergence rate. Notably, LoRA needs 20 epochs to achieve the same performance level that OFT attains after only 8 epochs. Importantly, the number of trainable parameters for OFT and COFT is comparable to LoRA (all significantly lower than ControlNet), yet their convergence is markedly more efficient. Additionally, the swift convergence exhibited by OFT is corroborated by the outcomes in Figure~\ref{fig:teaser}; in particular, the right-side example illustrates that OFT attains effective convergence using only 5\% of the ADE20K dataset, demonstrating much greater data efficiency compared to ControlNet and LoRA. For our qualitative evaluation, we primarily compare OFT, COFT, and ControlNet, as ControlNet's landmark error is closest to ours. Figure~\ref{fig:face_convergence_qual} presents evidence that both OFT and COFT demonstrate steady convergence, with their synthesized face poses progressively matching the control landmarks. In contrast, ControlNet exhibits a sudden jump in controllability only after the 8th epoch—a phenomenon consistent with the abrupt convergence described in \cite{zhang2023adding}. The greater convergence stability seen in OFT ensures a smoother and more predictable training trajectory, which simplifies the process of tuning OFT hyperparameters for practical applications.

\textbf{Quantitative comparison}. To assess controllable generation, we propose a control consistency metric that quantitatively measures alignment between the control signal derived from the generated image and the original input signal. Specifically, for the C2I task, we use the IoU and F1 score as evaluation criteria. For the S2I task, we utilize mean IoU, along with both mean and overall accuracy. The L2F task is assessed by calculating the mean $\ell_2$ distance between input and predicted landmark locations. Detailed explanations of these consistency metrics can be found in Appendix~\ref{app:settings}. Across all evaluated methods, optimal hyperparameters were selected. As summarized in Table~\ref{table:control_gen}, both OFT and COFT significantly outperform competing techniques in terms of control precision and robustness. Adapter-based models (such as T2I-Adapter and ControlNet) demonstrate slower convergence and ultimately underperform in final metrics. When compared with ControlNet, LoRA achieves superior performance on the S2I task, but lags behind in C2I and L2F. Overall, our experiments suggest that LoRA exhibits a limited upper bound in performance, even after extensive hyperparameter optimization. Conversely, OFT demonstrates ongoing performance gains as its number of trainable parameters is increased, indicating unsaturated improvement. Notably, our approach to quantitatively benchmarking controllable generation constitutes an original contribution, as it enables reliable measurement of fine-tuned models’ control fidelity—limited only by the accuracy of the underlying segmentation or detection module.

\textbf{Qualitative comparison}. In addition to quantitative metrics, we provide a qualitative assessment of OFT and COFT against leading approaches such as ControlNet, T2I-Adapter, and LoRA. As illustrated by the randomly sampled outputs in Figure~\ref{fig:control_final}, both OFT and COFT consistently deliver images of superior fidelity and realism while enabling precise control over generation. In the S2I task, LoRA fails to adhere to the given segmentation map, whereas ControlNet, OFT, and COFT effectively respect the input structure. Notably, compared to ControlNet, our methods (OFT and COFT) generate visually richer and more geometrically plausible images with substantially fewer parameters. For the C2I task, OFT and COFT succeed in synthesizing semantically coherent images from coarse Canny edges, an area where T2I-Adapter and LoRA lag behind. Regarding the L2F task, our approach attains the most precise pose alignment, even for faces exhibiting difficult poses. Across all three control tasks, OFT and COFT outperform existing baselines in visual quality, underscoring the efficacy of the OFT framework for controllable image generation. Further qualitative results for all three tasks can be found in Appendix~\ref{app:control_exp}. Additionally, Appendix~\ref{app:more_control_task} provides evidence that OFT generalizes well to other control scenarios such as dense pose to human body, sketch-to-image, and depth-to-image tasks.

\section{Concluding Remarks and Open Problems}

Inspired by the key role that angular relationships among neurons play in shaping visual semantics, this work introduces a straightforward yet powerful finetuning approach—orthogonal finetuning—to enhance control over text-to-image diffusion models. Our focus is on two primary applications: subject-driven generation and controllable generation. In contrast to prior methods, OFT achieves superior controllability and stability during finetuning while reducing the number of trainable parameters. Crucially, OFT also maintains the model’s inference efficiency without adding computational overhead, thereby supporting practical deployment.

This work with OFT gives rise to several compelling open questions. Firstly, OFT relies on Cayley parametrization to enforce orthogonality, which requires a matrix inverse operation and, as such, imposes some constraints on scalability. While we employ block diagonal parametrization as a partial solution, identifying a faster, differentiable means of computing the matrix inverse remains unresolved. Secondly, OFT inherently supports compositionality: orthogonal matrices generated by distinct OFT finetuning processes can be multiplied together to yield another orthogonal matrix. A pressing open question is whether the combined set of orthogonal matrices effectively retains the learned information from all corresponding downstream tasks. Lastly, OFT’s parameter efficiency is closely tied to its block diagonal parametrization, a design that inevitably adds bias and restricts flexibility. Seeking strategies to further enhance parameter efficiency while minimizing bias represents a significant avenue for future investigation.

\end{document}